% ============================================================
% Chapter 5: Development
%   Section: News Acquisition: Discovery and Extraction
%     Subsection: Overview
% Included from chapters/05-development/02-news-acquisition/news-acquisition.tex
% ============================================================

\subsection{Overview}
\label{subsec:scraping-overview}

Acquisition is two steps, composed by the \texttt{Scraper} class
(Figure~\ref{fig:acquisition}):

\begin{enumerate}
    \item \textbf{Discovery} (\texttt{DiscoveryService}) takes a
    configured source and returns the URLs of its recent articles,
    together with whatever the source says about each one before it is
    opened: a title, a summary and a publication time.
    \item \textbf{Extraction} (\texttt{ExtractorService}) takes one
    URL and returns the article: its body, title, author, date and,
    when the page states it, the publication time to the minute.
\end{enumerate}

The two steps are used separately. Discovery feeds the candidate
rounds from which a person or the AI chooses what to analyse
(Section~\ref{sec:design-selection}), the one-shot ingestion endpoint
and the daily labelling batch. Extraction is called for every chosen
candidate, for every URL posted to the analysis endpoint, for every
web page read as evidence for a claim, and by the tools that check the
sources' health. Each call states its \emph{purpose}
(\texttt{article}, \texttt{evidence}, \texttt{enrichment},
\texttt{ingestion}, \texttt{discovery}, \texttt{source\_check},
\texttt{probe}, \texttt{labelling}), which decides how much the call
may spend and under which heading it is counted.

Four principles run through the implementation, each introduced after
a measured failure that is cited where it is described:

\begin{itemize}
    \item \textbf{Cheapest first.} Discovery and extraction both try
    their strategies in order of cost, and only go to the next one when
    it could plausibly succeed where the previous one failed.
    \item \textbf{Every fetch is guarded.} Every URL is checked against
    an address-based guard before any request is sent, and again at
    every redirect.
    \item \textbf{Everything is counted.} Every page wanted is recorded
    with its outcome, the strategies tried and the requests actually
    sent, so a source that stops working shows up as such rather than
    as an absence of news.
    \item \textbf{Sources are configuration, not code.} Adding a source
    means adding a YAML file; nothing in the code names a particular
    outlet.
\end{itemize}

\begin{figure}[ht]
\centering
\begin{tikzpicture}[
    node distance=6mm and 14mm,
    box/.style={draw, rounded corners=2pt, align=center,
                font=\footnotesize, text width=5.2cm,
                minimum height=7mm, inner sep=3pt},
    io/.style={box, fill=black!6},
    stage/.style={box, fill=ganttplan!20},
    gate/.style={box, fill=ganttextra!30},
    arr/.style={-{Stealth[length=2mm]}, thick},
    alt/.style={-{Stealth[length=2mm]}, dashed},
    lbl/.style={font=\scriptsize\itshape, fill=white, inner sep=1pt},
]
% ---- Discovery ----
\node[io]                   (src)    {Source catalogue\\(one YAML file per source)};
\node[stage, below=of src]  (rss)    {1. The source's RSS feed};
\node[stage, below=of rss]  (feeds)  {2. Feed discovery (trafilatura)};
\node[stage, below=of feeds](pages)  {3. Topic section pages};
\node[gate,  below=of pages](filt)   {Discovery filters: shape, topic,\\
                                      off-mission sections, age,\\
                                      already stored};
\node[io,    below=of filt] (round)  {Candidates $\rightarrow$ selection};

\draw[arr] (src) -- (rss);
\draw[alt] (rss) -- node[lbl]{nothing found} (feeds);
\draw[alt] (feeds) -- node[lbl]{nothing found} (pages);
\draw[arr] (pages) -- (filt);
\draw[arr] (filt) -- (round);
\coordinate (dl) at ($(rss.west)+(-5mm,0)$);
\draw[thick] (rss.west) -- (dl);
\draw[thick] (feeds.west) -- (dl |- feeds.west);
\draw[arr] (dl) |- (filt.west);

% ---- Extraction ----
\node[io,    right=of src]  (in)     {A URL: a chosen candidate, a posted\\article, an evidence page};
\node[gate,  below=of in]   (guard)  {URL guard, one shared fetch};
\node[stage, below=of guard](traf)   {1. trafilatura};
\node[stage, below=of traf] (bs)     {2. BeautifulSoup, same HTML};
\node[stage, below=of bs]   (pw)     {3. Headless browser (Playwright)};
\node[io,    below=of pw]   (out)    {Validated article and metadata};

\draw[arr] (in) -- (guard);
\draw[arr] (guard) -- (traf);
\draw[alt] (traf) -- node[lbl]{no article in the page} (bs);
\draw[alt] (bs) -- node[lbl]{still none} (pw);
\draw[alt] (traf.east) -- ++(5mm,0) |- node[lbl, pos=0.25, right, xshift=1mm]{401/403/429} (pw.east);
\draw[arr] (pw) -- (out);
\coordinate (el) at ($(traf.west)+(-5mm,0)$);
\draw[thick] (traf.west) -- (el);
\draw[thick] (bs.west) -- (el |- bs.west);
\draw[arr] (el) |- (out.west);

\node[font=\footnotesize\bfseries, above=2mm of src] {Discovery};
\node[font=\footnotesize\bfseries, above=2mm of in]  {Extraction};
\end{tikzpicture}
\caption{The acquisition layer. Solid arrows are the normal path;
dashed arrows are escalations, taken only when the previous step could
not help. A missing page, a timeout or a guard refusal stops the
extraction at once.}
\label{fig:acquisition}
\end{figure}
